\section*{Capítulo \ref{ch:g11:binom} --- A distribuição binomial}

\begin{solution}{exo:g11:binom:1}
Linhas $6$ e $7$:
\[
1,\ 6,\ 15,\ 20,\ 15,\ 6,\ 1
\qquad\text{e}\qquad
1,\ 7,\ 21,\ 35,\ 35,\ 21,\ 7,\ 1 .
\]
Logo $\binom62 = 15$, $\binom73 = 35$, $\binom74 = 35$ (a simetria
$\binom73 = \binom74$ reflete a troca entre sucessos e fracassos).
\end{solution}

\begin{solution}{exo:g11:binom:2}
Número fixo de ensaios ($4$ lançamentos), dois resultados por lançamento
(seis ou não, $p = \frac16$) e lançamentos independentes:
$X \sim \mathcal B(4, \frac16)$.
\[
\P(X=0) = \left(\frac56\right)^4 = \frac{625}{1296} \approx 0.48,
\qquad
\P(X=1) = 4 \times \frac16\left(\frac56\right)^3 = \frac{500}{1296}
\approx 0.39,
\]
\[
\P(X \geq 2) = 1 - \frac{625 + 500}{1296} = \frac{171}{1296}
\approx 0.13 .
\]
\end{solution}

\begin{solution}{exo:g11:binom:3}
\emph{1.} Binomial $\mathcal B(20, \frac12)$: $n$ fixo, mesmo $p$,
lançamentos independentes.

\emph{2.} Não é binomial: as cartas são distribuídas \emph{sem
reposição}, de modo que a \omterm{def:g10:proba:distribution}{probabilidade} de um ás muda de carta para carta
e as retiradas não são independentes.

\emph{3.} Binomial $\mathcal B(7, 0.3)$, pela hipótese de independência
enunciada.
\end{solution}

\begin{solution}{exo:g11:binom:4}
$\E(X) = 50 \times 0.2 = 10$;
$\V(X) = 50 \times 0.2 \times 0.8 = 8$;
$\sigma(X) = 2\sqrt2 \approx 2.83$.
\end{solution}

\begin{solution}{exo:g11:binom:5}
$X \sim \mathcal B(6, 0.7)$.
\[
\P(X = 4) = \binom64 (0.7)^4 (0.3)^2 = 15 \times 0.2401 \times 0.09
\approx 0.324 .
\]
\[
\P(X \geq 5) = \binom65 (0.7)^5(0.3) + (0.7)^6
= 6 \times 0.16807 \times 0.3 + 0.117649 \approx 0.420 .
\]
\end{solution}

\begin{solution}{exo:g11:binom:6}
$X \sim \mathcal B\left(8, \frac12\right)$; todo caminho tem
\omterm{def:g10:proba:distribution}{probabilidade} $\frac{1}{256}$, logo
\[
\P(X \geq 6) = \frac{\binom86 + \binom87 + \binom88}{256}
= \frac{28 + 8 + 1}{256} = \frac{37}{256} \approx 0.14 .
\]
Chutar aprova cerca de uma vez em sete.
\end{solution}

\begin{solution}{exo:g11:binom:7}
O \omterm{def:g10:proba:operations}{complementar} de ``pelo menos uma de cada'' é ``as cinco iguais'':
todas A ou todas B, cada caso com \omterm{def:g10:proba:distribution}{probabilidade}
$\left(\frac12\right)^5 = \frac1{32}$. Logo
\[
\P(\text{uma de cada tipo}) = 1 - \frac{2}{32} = \frac{15}{16} .
\]
\end{solution}

\begin{solution}{exo:g11:binom:8}
$\P(X = 3) = p^3$ e $\P(X \geq 1) = 1 - (1-p)^3$. Resolver
$p^3 = \frac{27}{64} = \left(\frac34\right)^3$ dá $p = \frac34$ (a \omterm{def:g11:func:function}{função}
cubo é estritamente \omterm{def:g11:func:monotone}{crescente}, \cref{ch:g11:func}, de modo que a solução
é única).
\end{solution}

\begin{solution}{exo:g11:binom:9}
$\P(\text{pelo menos uma cara}) = 1 - \left(\frac12\right)^n$, de modo
que a condição é $\left(\frac12\right)^n < 0.01$, isto é, $2^n > 100$.
Como $2^6 = 64$ e $2^7 = 128$: a partir de $n = 7$ lançamentos.
\end{solution}

\begin{solution}{exo:g11:binom:10}
A soma da linha $n$ conta todos os caminhos da \omterm{met:g10:proba:tree}{árvore} de $n$ ensaios,
separados pelo número de sucessos. Mas a \omterm{met:g10:proba:tree}{árvore} dobra seus caminhos a
cada ensaio (cada caminho se divide em S e F), de modo que ela tem $2^n$
caminhos ao todo. Logo $\sum_{k} \binom nk = 2^n$. Alternativamente, por
indução: a linha $0$ soma $1 = 2^0$, e a relação de Pascal faz cada
entrada da linha $n$ contribuir para exatamente duas entradas da linha
$n+1$, de modo que as somas das linhas dobram.
\end{solution}

\begin{solution}{exo:g11:binom:11}
\emph{1.} Sob a afirmação, $X \sim \mathcal B(10, 0.6)$. Somando os
primeiros termos:
\begin{align*}
\P(X \leq 3)
&= (0.4)^{10} + 10(0.6)(0.4)^9 + 45(0.6)^2(0.4)^8
+ 120(0.6)^3(0.4)^7\\
&\approx 0.0001 + 0.0016 + 0.0106 + 0.0425 = 0.0548 .
\end{align*}

\emph{2.} Um resultado pelo menos tão extremo quanto o observado ($3$
aprovações ou menos) tem \omterm{def:g10:proba:distribution}{probabilidade} de cerca de $5.5\%$ --- logo
\emph{acima} do limiar de $5\%$. Aplicando a regra ao pé da letra, a
observação é (por pouco) compatível com a afirmação e não se rejeita a
hipótese. O exemplo mostra o quanto as decisões no limite são sensíveis à
escolha do limiar: com uma convenção de $6\%$ a conclusão se inverteria.
\end{solution}

\begin{solution}{pb:g11:binom:1}
\textbf{1.} Linhas: $1$; $1\,1$; $1\,2\,1$; $1\,3\,3\,1$;
$1\,4\,6\,4\,1$; $1\,5\,10\,10\,5\,1$;
$1\,6\,15\,20\,15\,6\,1$. Simetria: escolher quais $k$ objetos
levar é o mesmo ato que escolher quais $n - k$ deixar para
trás.

\textbf{2.} $1 + 4 + 6 + 4 + 1 = 16 = 2^4$;
$1 + 5 + 10 + 10 + 5 + 1 = 32 = 2^5$. Demonstração: os
$\binom nk$ contam os subconjuntos de cada tamanho de um
conjunto de $n$ elementos, e os subconjuntos ao todo são $2^n$
(cada elemento entra ou não, de forma independente).

\textbf{3.} Comissões de $k$ pessoas escolhidas entre $n + 1$,
uma das quais é Zoé: as sem Zoé são $\binom nk$ (escolher as
$k$ entre as outras); as com Zoé são $\binom{n}{k-1}$
(escolher os $k - 1$ colegas dela). Total:
$\binom nk + \binom{n}{k-1}$.

\textbf{4.} Linha $7$ do triângulo: $1\,7\,21\,35\,\dots$: $35$.
Fórmula: $\frac{7 \times 6 \times 5}{3 \times 2 \times 1} =
35$.

\textbf{5.} $1 + 3 + 6 + 10 = 20 = \binom63$. Cascata:
$\binom63 = \binom52 + \binom53 = \binom52 + \binom42 +
\binom43 = \binom52 + \binom42 + \binom32 + \binom33$ --- cada
aplicação da relação de Pascal descasca um degrau da escada.

\textbf{6.} Número fixo $n$ de quiques; cada quique um \omterm{def:g11:binom:bernoulli}{ensaio de
Bernoulli} independente com o mesmo $p = \frac12$; o número da
canaleta conta os sucessos (quiques para a direita): as três
caixas do \cref{met:g11:binom:recognize} marcadas:
$\mathcal B\!\left(n, \frac12\right)$.

\textbf{7.} \omterm{def:g10:proba:distribution}{Probabilidades} $\frac{1}{16}, \frac{4}{16},
\frac{6}{16}, \frac{4}{16}, \frac{1}{16}$ para as canaletas
$0, \dots, 4$: a canaleta central $2$ é a mais cheia.

\textbf{8.} Quantidades esperadas $= 1024 \times
\binom{10}{k}/1024 = \binom{10}{k}$: canaleta central,
$\binom{10}{5} = 252$ bolinhas; canaleta $7$:
$\binom{10}{7} = 120$; cada canaleta da ponta: $1$ bolinha. Um
centro alto que decai simetricamente até pontas finíssimas: o
sino.

\textbf{9.} $\E = np = 5$; $V = np(1 - p) = 2.5$;
$\sigma \approx 1.58$. A menos de $2\sigma$: as canaletas $2$ a
$8$ carregam
\[
\frac{45 + 120 + 210 + 252 + 210 + 120 + 45}{1024}
= \frac{1002}{1024} \approx 98\,\%
\]
das bolinhas --- muito melhor do que os $75\,\%$ genéricos de
Chebyshev (\cref{pb:g11:stat:1}): as formas de sino se
concentram com força.

\textbf{10.} $\E = 6$, $V = 10 \times 0.6 \times 0.4 = 2.4$,
$\sigma \approx 1.55$: a pilha mantém a forma de sino, mas
desliza o cume para a canaleta $6$ --- um tabuleiro inclinado é
uma moeda viciada tornada visível.

\textbf{11.} O número da canaleta é uma \emph{soma} de muitos
empurrões aleatórios pequenos, independentes e do mesmo
tamanho --- e somas assim sempre se organizam em sino: a
maioria dos empurrões se cancela, e os extremos exigem
unanimidade. Alturas, erros de medição e inúmeras grandezas
naturais são igualmente somas de muitos efeitos pequenos e
independentes, e é por isso que a mesma silhueta aparece em
toda parte; o teorema que certifica isso é o teorema central do
limite.

\textbf{12.} Uma varrida: um time vence os $4$:
$2 \times \left(\frac12\right)^4 = \frac18$.

\textbf{13.} Sete jogos exigem $3$--$3$ depois de seis:
$\binom63 \left(\frac12\right)^6 = \frac{20}{64} =
\frac{5}{16}$.

\textbf{14.} Termina em $5$: o vencedor leva o jogo 5 e $3$ dos
$4$ primeiros: $2 \times \binom43 \left(\frac12\right)^5 =
\frac14$. Termina em $6$:
$2 \times \binom53 \left(\frac12\right)^6 = \frac{5}{16}$.
\omterm{def:g11:prob:rv}{Distribuição} sobre $4, 5, 6, 7$:
$\frac18, \frac14, \frac{5}{16}, \frac{5}{16}$ (soma $1$).
Duração esperada: $4 \cdot \frac18 + 5 \cdot \frac14 + 6 \cdot
\frac{5}{16} + 7 \cdot \frac{5}{16} = 5.8125$ jogos. As séries
de seis e sete jogos são as mais prováveis --- o drama já vem
embutido no formato.

\textbf{15.} Vitória em $4$: $0.6^4 = 0.1296$; em $5$:
$\binom43\,0.6^3 \times 0.4 \times 0.6 = 0.2074$; em $6$:
$\binom53\,0.6^3 \times 0.4^2 \times 0.6 = 0.2074$; em $7$:
$\binom63\,0.6^3 \times 0.4^3 \times 0.6 = 0.1659$. Total:
cerca de $0.710$: um time de $60\,\%$ por jogo vence $71\,\%$
das séries --- a série amplifica a vantagem.

\textbf{16.} Final única: $60\,\%$. Melhor de 3:
$p^2 + 2p^2 q = 0.36 + 0.288 = 0.648$. A escada
$60\,\% \to 65\,\% \to 71\,\%$ continua com o comprimento: mais
jogos tiram a \omterm{def:g11:stat:mean}{média} da sorte (a lei dos grandes números em
miniatura), de modo que finais longas coroam a habilidade ---
que é precisamente o que as ligas vendem.

\textbf{17.} Moeda honesta: $\E = 50$,
$\sigma = \sqrt{25} = 5$; $z = \frac{62 - 50}{5} = 2.4$: além
da convenção de $2\sigma$ --- a moeda merece uma investigação.

\textbf{18.} $\P(X = 0) = 0.98^{50} \approx 0.364$;
$\P(X = 1) = 50 \times 0.02 \times 0.98^{49} \approx 0.372$;
$\P(X = 2) = \binom{50}{2} 0.02^2 \times 0.98^{48} \approx
0.186$. Logo $\P(X \geq 3) \approx 1 - 0.922 = 0.078$: cerca de
$7.8\,\%$ --- acima do limiar de $5\,\%$, portanto ainda não
basta para rejeitar a afirmação; um segundo lote ruim mudaria a
história.

\textbf{19.} $\P(\text{pelo menos uma vitória}) = 1 -
0.999^{1000} \approx 1 - 0.368 = 0.632$: mil bilhetes a um em
mil não dão certeza, e sim $63\,\%$. O recorrente $0.368$ é
$\frac1e$ disfarçado --- a constante $e$ faz sua entrada oficial
no ano 12.

\textbf{20.} Reconhecimento: $n$ fixo, independência, $p$
constante --- só então, binomial. Tabela: o \omterm{prop:g11:binom:pascal}{triângulo de
Pascal}, linha $n$. Forma: o sino, simétrico para $p = \frac12$
e deslocado nos demais casos. Bússola: centro $np$, dispersão
$\sqrt{np(1-p)}$ --- os \omterm{pb:g11:stat:1}{escores z} das decisões. Herdeiros: a
curva em sino contínua de que as pilhas se aproximam e a lei
dos grandes números, que explica por que tabuleiros grandes
nunca mentem.
\end{solution}
